% !TeX root = closed-systems-from-comparison-completeness.tex
% ============================================================
\chapter{Transport-Closed Dynamics and Spacetime Interleaving}
\label{ch:stitching}
% ============================================================
\section{Introduction}
Under the standing principle of closed-world admissibility
(\cref{sp:closed-world}), this chapter determines the transport-closed
dynamics consequences of \cref{sp:quotient,sp:transport} on the closed
side of the stack.  Using the observer-side irreversible filtration
inherited from \cref{ch:flow}, it then assembles the canonical stitched
refinement/observer arena \(\mathsf{ST}\) used immediately in
\cref{ch:christoffel} and later in \cref{ch:causality,ch:connection}.
\Cref{ch:flow} isolated the source of irreversibility in the
closed-system stack.  The intrinsic quotient
\[
	\pi\colon X\to \Phys:=X/G
\]
does not itself generate an arrow of time: reversible evolution on \(X\)
descends to reversible evolution on \(\Phys\).  Irreversibility enters
only at the later internal observer-induced descent on \(\Phys\), and
the observer-side temporal filtration used here is inherited from that
stage.
The purpose of the present chapter is different.  We return to the
closed side of the framework and determine what transport closure entails
on admissible dynamics before any observer reduction is imposed.
Two facts are proved.
\begin{enumerate}[label=\textup{(\roman*)}]
	\item
	      Under closed-system semantics and transport closure, every admissible
	      dynamical evolution preserves the primitive comparison structure and
	      therefore lies in the intrinsic symmetry group
	      \[
		      G=\Aut(U,\mathcal C).
	      \]
	\item
	      Once comparison preservation is determined, admissible dynamics acts
	      canonically on the refinement tower of Boolean algebras and hence on
	      its Stone inverse limit.  If an observer reduction on the limit
	      descends to finite refinement levels, then temporal coarse-graining and
	      spatial refinement interleave functorially, yielding a canonical
	      two-parameter inverse system.
\end{enumerate}
To state the second conclusion in an intrinsically categorical form, we
prove not merely that each admissible symmetry acts on each refinement
stage, but that
\[
	\Aut(U,\mathcal C)
\]
acts by automorphisms of the entire refinement inverse system.
The resulting two-parameter limit object is the first object in the
stack in which refinement and observer-induced temporal coarsening are
assembled into a single canonical construction.
\medskip
The logical position of the chapter is therefore
\[
	\begin{aligned}
		 & \text{transport closure}
		\Longrightarrow
		\text{comparison-preserving dynamics}
		\Longrightarrow
		\text{refinement action}
		\\
		 & \Longrightarrow
		\text{refinement/observer interleaving}
		\Longrightarrow
		\mathsf{ST}.
	\end{aligned}
\]
No curvature, smooth realization, or macroscopic field law appears
here.  Those belong to later chapters.
% ------------------------------------------------------------
\section{Transport closure determines comparison preservation}
% ------------------------------------------------------------
We begin by fixing the meaning of transport closure at the present
level of the stack.
\begin{definition}[Transport-closed extension]
	An extension of quotient semantics is called
	\emph{transport-closed} if it introduces neither
	\begin{enumerate}[label=\textup{(\roman*)}]
		\item object-locus representative data, i.e.\ no section
		      \[
			      s\colon \Phys\to X
			      \qquad\text{with}\qquad
			      \pi\circ s=\id_{\Phys},
		      \]
		      nor
		\item morphism-locus obstruction, i.e.\ no nontrivial holonomy or loop
		      defect.
	\end{enumerate}
	Equivalently, by the descent--obstruction classification together with
	the loop criterion from the transport spine
	(\cref{thm:structural-closure} and the endpoint-determined criterion in
	the descent--obstruction chapter), admissible transport is
	endpoint-determined.
\end{definition}
\begin{definition}[Admissible dynamics]
	A dynamical evolution on \(U\) is called \emph{admissible} if it is
	realized within a transport-closed extension of quotient semantics.
\end{definition}
The first result removes an external hypothesis.  Comparison
preservation is not assumed; it is determined.
\begin{lemma}[Transport-closed endpoint-determined realizations add no new comparison data]
	\label{lem:endpoint-no-new-comparison}
	Under the standing principle \cref{sp:closed-world}, assume
	closed-system semantics and transport closure.  Any
	endpoint-determined realization of quotient-level dynamics introduces
	no comparison distinction beyond the descended content on \(\Phys\).
\end{lemma}
\begin{proof}
	By the two-locus classification, any admissible enrichment beyond
	quotient semantics is exhausted by object-locus representative choice
	and morphism-locus transport obstruction, possibly with both and with
	no third locus.
	Endpoint-determined transport excludes the morphism locus, while
	transport closure excludes the object locus.  Hence no additional
	comparison distinction can be introduced beyond the quotient-level
	content already determined by \(\Phys\).
\end{proof}
\begin{theorem}[Transport closure determines comparison preservation]
	\label{thm:transport-forces-Aut}
	Under the standing principle \cref{sp:closed-world}, assume
	closed-system semantics and transport closure.  Then every
	admissible dynamical transformation on \(U\) preserves every comparison
	predicate and therefore belongs to the intrinsic symmetry group
	\[
		G=\Aut(U,\mathcal C).
	\]
\end{theorem}
\begin{proof}
	Let \(\Phi\colon U\to U\) be admissible.  Suppose
	\[
		\Phi\notin \Aut(U,\mathcal C).
	\]
	Then, by definition of \(\Aut(U,\mathcal C)\), there exist
	\(c\in\mathcal C\) and \(u,v\in U\) such that
	\[
		c(\Phi(u),\Phi(v))\neq c(u,v).
	\]
	Under rectangular completeness, the comparison world admits its
	canonical diagonal presentation
	\[
		U\cong X:=X_A\times X_B,
		\qquad
		\pi\colon X\to \Phys:=X/G,
	\]
	and closed-system semantics identifies admissible state content with
	data descending through \(\pi\).  Thus any admissible realization of
	\(\Phi\) must respect the quotient boundary determined by \(\pi\).
	But the inequality
	\[
		c(\Phi(u),\Phi(v))\neq c(u,v)
	\]
	shows that \(\Phi\) produces a comparison distinction not already
	determined by the descended quotient content on \(\Phys\).  By
	\cref{lem:endpoint-no-new-comparison}, no endpoint-determined
	realization can produce such a distinction.  Hence any realization of
	\(\Phi\) must involve enrichment beyond endpoint-determined quotient
	semantics.
	By the universal two-locus classification of enrichment
	(\cref{thm:classification} in \cref{ch:rotation}, equivalently the
	action-groupoid classification together with the section--holonomy
	dichotomy), every such non-endpoint-determined realization is
	exhausted by representative choice and morphism-level transport,
	possibly with both and with no third locus.
	Transport closure excludes both.  This contradiction proves that
	\(\Phi\in\Aut(U,\mathcal C)\).
\end{proof}
\begin{remark}[Derived symmetry]
	Theorem~\ref{thm:transport-forces-Aut} replaces the external hypothesis
	that admissible dynamics preserves comparisons.  In the closed-system
	stack, comparison preservation is not an independent postulate.  It is
	a consequence of admissibility together with transport closure.
\end{remark}
% ------------------------------------------------------------
\section{Refinement tower and Stone inverse limit}
% ------------------------------------------------------------
We now pass from the comparison layer to the refinement layer.
Let
\[
	(\B_k)_{k\in\mathbb N}
\]
be the refinement tower of \cref{ch:foundation}, and let
\[
	\B_\infty
	:=
	\left\langle \bigcup_{k\in\mathbb N}\B_k\right\rangle
\]
be the generated limit Boolean algebra.  Write
\[
	S_k:=\UF(\B_k),
	\qquad
	S_\infty:=\UF(\B_\infty)
\]
for the corresponding Stone spaces.  Restriction of ultrafilters along
the inclusions
\[
	\B_k\subseteq\B_{k+1}
\]
induces bonding maps
\[
	\phi_{k+1,k}\colon S_{k+1}\to S_k.
\]
\begin{proposition}[Stone inverse limit]
	\label{prop:Stone-limit}
	There is a canonical homeomorphism
	\[
		S_\infty\cong \varprojlim_k S_k.
	\]
	We denote the inverse-limit projections by
	\[
		p_k\colon S_\infty\to S_k.
	\]
\end{proposition}
\begin{proof}
	This is the inverse-limit representation of admissible ultrafilters
	proved in the foundational chapter, applied to the present refinement
	tower.
\end{proof}
\begin{lemma}[Surjectivity of the projections]
	\label{lem:pk-surjective}
	For every \(k\), the projection
	\[
		p_k\colon S_\infty\to S_k
	\]
	is surjective.
\end{lemma}
\begin{proof}
	Each bonding map
	\[
		\phi_{k+1,k}\colon S_{k+1}\to S_k
	\]
	is induced by restriction of ultrafilters along
	\[
		\B_k\subseteq \B_{k+1}.
	\]
	Every ultrafilter on \(\B_k\) extends to an ultrafilter on
	\(\B_{k+1}\), so each \(\phi_{k+1,k}\) is surjective.  The coordinate
	projections from an inverse limit of nonempty compact Hausdorff spaces
	with surjective bonding maps are therefore surjective.  Under the
	identification of \cref{prop:Stone-limit}, this is exactly the
	surjectivity of \(p_k\).
\end{proof}
% ------------------------------------------------------------
\section{Functorial action on the refinement inverse system}
% ------------------------------------------------------------
Once admissible dynamics is determined into \(\Aut(U,\mathcal C)\), it acts
not only on \(U\) but on every comparison-generated Boolean stage, and
hence on the full Stone tower.
\begin{proposition}[Boolean invariance]
	\label{prop:boolean-inv}
	Every comparison-preserving bijection
	\[
		\phi\in\Aut(U,\mathcal C)
	\]
	acts on \(\B_\infty\) by pullback
	\[
		E\longmapsto \phi^{-1}(E),
	\]
	preserves each finite-stage algebra \(\B_k\), and therefore induces
	maps
	\[
		\widehat\phi^{(k)}\colon S_k\to S_k,
		\qquad
		\widehat\phi^\infty\colon S_\infty\to S_\infty.
	\]
\end{proposition}
\begin{proof}
	It is enough to check preservation on the generating comparison sets.
	For a basic left comparison set
	\[
		L_{c,w}:=\{u\in U:c(u,w)=1\},
	\]
	we have
	\[
		u\in\phi^{-1}(L_{c,w})
		\iff
		c(\phi(u),w)=1
		\iff
		c(u,\phi^{-1}(w))=1
		\iff
		u\in L_{c,\phi^{-1}(w)}.
	\]
	Thus pullback sends left generators to left generators.
	Similarly, for a basic right comparison set
	\[
		R_{c,w}:=\{u\in U:c(w,u)=1\},
	\]
	one computes
	\[
		u\in\phi^{-1}(R_{c,w})
		\iff
		c(w,\phi(u))=1
		\iff
		c(\phi^{-1}(w),u)=1
		\iff
		u\in R_{c,\phi^{-1}(w)}.
	\]
	Hence right generators are preserved as well.
	Therefore pullback preserves the generating comparison family, hence
	every \(\B_k\) and therefore \(\B_\infty\).  The induced maps on Stone
	spaces are the stated maps.
\end{proof}
\begin{theorem}[Functorial action on the refinement inverse system]
	\label{thm:Aut-functorial}
	The assignment
	\[
		\phi\longmapsto\bigl(\widehat\phi^{(k)}\bigr)_{k\in\mathbb N}
	\]
	defines an action of
	\[
		\Aut(U,\mathcal C)
	\]
	by automorphisms of the inverse system
	\[
		\bigl(S_k,\phi_{k+1,k}\bigr)_{k\in\mathbb N}.
	\]
	Equivalently, for every
	\[
		\phi\in\Aut(U,\mathcal C)
	\]
	and every \(k\), the square
	\[
		\begin{tikzcd}
			S_{k+1} \arrow[r,"\widehat\phi^{(k+1)}"] \arrow[d,"\phi_{k+1,k}"']
			& S_{k+1} \arrow[d,"\phi_{k+1,k}"] \\
			S_k \arrow[r,"\widehat\phi^{(k)}"']
			& S_k
		\end{tikzcd}
	\]
	commutes, and
	\[
		\widehat{\id}^{(k)}=\id_{S_k},
		\qquad
		\widehat{(\phi\circ\psi)}^{(k)}
		=
		\widehat\phi^{(k)}\circ\widehat\psi^{(k)}
	\]
	for all \(k\).
\end{theorem}
\begin{proof}
	Fix \(\phi\in\Aut(U,\mathcal C)\).  By
	\cref{prop:boolean-inv}, the maps
	\[
		\widehat\phi^{(k)}\colon S_k\to S_k
	\]
	are well-defined.
	To verify compatibility with the bonding maps, let
	\(\mathfrak u\in S_{k+1}\).  Then
	\[
		\phi_{k+1,k}\bigl(\widehat\phi^{(k+1)}(\mathfrak u)\bigr)
		=
		\widehat\phi^{(k+1)}(\mathfrak u)\cap\B_k.
	\]
	By definition,
	\[
		E\in\widehat\phi^{(k+1)}(\mathfrak u)
		\iff
		\phi^{-1}(E)\in\mathfrak u
		\qquad
		(E\in\B_{k+1}).
	\]
	Restricting to \(E\in\B_k\), we obtain
	\[
		E\in\phi_{k+1,k}\bigl(\widehat\phi^{(k+1)}(\mathfrak u)\bigr)
		\iff
		\phi^{-1}(E)\in \mathfrak u\cap\B_k
		\iff
		E\in\widehat\phi^{(k)}(\mathfrak u\cap\B_k).
	\]
	Thus
	\[
		\phi_{k+1,k}\circ\widehat\phi^{(k+1)}
		=
		\widehat\phi^{(k)}\circ\phi_{k+1,k}.
	\]
	Identity and composition are immediate from the corresponding facts for
	pullback on subsets of \(U\).  The identity automorphism induces the
	identity on every Boolean algebra, hence
	\[
		\widehat{\id}^{(k)}=\id_{S_k}.
	\]
	If \(\phi,\psi\in\Aut(U,\mathcal C)\), then
	\[
		(\phi\circ\psi)^{-1}=\psi^{-1}\circ\phi^{-1},
	\]
	so
	\[
		\widehat{(\phi\circ\psi)}^{(k)}
		=
		\widehat\phi^{(k)}\circ\widehat\psi^{(k)}.
	\]
	Therefore \(\Aut(U,\mathcal C)\) acts by automorphisms of the inverse
	system.
\end{proof}
\begin{corollary}[Action on the Stone limit]
	\label{cor:Aut-limit-action}
	The action of \(\Aut(U,\mathcal C)\) on the inverse system induces an
	action on the inverse limit \(S_\infty\).  Under the identification
	\[
		S_\infty\cong \varprojlim_k S_k,
	\]
	this action agrees with the Stone action
	\[
		\widehat\phi^\infty\colon S_\infty\to S_\infty.
	\]
\end{corollary}
\begin{proof}
	By \cref{thm:Aut-functorial}, each \(\phi\in\Aut(U,\mathcal C)\)
	determines a compatible family of maps on the inverse system and hence
	a unique induced map on the inverse limit.  By construction, this is
	the Stone-dual map induced directly from pullback on \(\B_\infty\).
\end{proof}
\begin{corollary}[Induced flow on the Stone limit]
	Let
	\[
		(\Phi_t)_{t\in\mathbb R}
	\]
	be an admissible reversible flow on \(U\).  Then each \(\Phi_t\)
	induces compatible homeomorphisms
	\[
		\widehat\Phi_t^{(k)}\colon S_k\to S_k,
		\qquad
		\widehat\Phi_t^\infty\colon S_\infty\to S_\infty,
	\]
	and
	\[
		(\widehat\Phi_t^\infty)_{t\in\mathbb R}
	\]
	is a reversible flow on \(S_\infty\).
\end{corollary}
\begin{proof}
	By \cref{thm:transport-forces-Aut}, each \(\Phi_t\) lies in
	\(\Aut(U,\mathcal C)\), so
	\cref{thm:Aut-functorial,cor:Aut-limit-action} produce the induced
	maps on every \(S_k\) and on \(S_\infty\).  Since
	\[
		\Phi_{t+s}=\Phi_t\circ\Phi_s,
		\qquad
		\Phi_0=\id_U,
	\]
	the induced maps satisfy
	\[
		\widehat\Phi_{t+s}^{(k)}
		=
		\widehat\Phi_t^{(k)}\circ\widehat\Phi_s^{(k)},
		\qquad
		\widehat\Phi_0^{(k)}=\id_{S_k},
	\]
	and likewise on \(S_\infty\).  Hence
	\((\widehat\Phi_t^\infty)_{t\in\mathbb R}\) is a reversible flow.
\end{proof}
% ------------------------------------------------------------
\section{Observer reduction on the Stone limit}
% ------------------------------------------------------------
We now reintroduce observer structure, but only at the limit level.
Let
\[
	\theta_\infty\colon S_\infty\twoheadrightarrow Y
\]
be a surjective observer reduction.
\begin{definition}[$k$-level admissibility]
	The observer reduction \(\theta_\infty\) is called
	\emph{\(k\)-level admissible} if it is constant on fibers of
	\[
		p_k\colon S_\infty\to S_k,
	\]
	that is,
	\[
		p_k(x)=p_k(y)
		\Longrightarrow
		\theta_\infty(x)=\theta_\infty(y).
	\]
\end{definition}
\begin{lemma}[Descent criterion]
	\label{lem:descent-criterion}
	Fix \(k\).  The following are equivalent:
	\begin{enumerate}[label=\textup{(\roman*)}]
		\item
		      There exists a map
		      \[
			      \theta_k\colon S_k\to Y
		      \]
		      such that
		      \[
			      \theta_k\circ p_k=\theta_\infty.
		      \]
		\item
		      \(\theta_\infty\) is \(k\)-level admissible.
	\end{enumerate}
	If such a map exists, it is unique.
\end{lemma}
\begin{proof}
	If \(\theta_k\) exists and \(p_k(x)=p_k(y)\), then
	\[
		\theta_\infty(x)=\theta_k(p_k(x))=\theta_k(p_k(y))=\theta_\infty(y),
	\]
	so \(\theta_\infty\) is constant on fibers of \(p_k\).
	Conversely, assume \(\theta_\infty\) is constant on fibers of \(p_k\).
	Given \(u\in S_k\), choose \(x\in S_\infty\) with \(p_k(x)=u\), which
	is possible by \cref{lem:pk-surjective}, and define
	\[
		\theta_k(u):=\theta_\infty(x).
	\]
	This is well-defined by fiber constancy.  The identity
	\[
		\theta_k\circ p_k=\theta_\infty
	\]
	is immediate, and uniqueness follows from the surjectivity of \(p_k\).
\end{proof}
\begin{definition}[Induced finite-level filtration]
	Assume \(\theta_\infty\) is \(k\)-level admissible, and let
	\[
		\theta_k\colon S_k\to Y
	\]
	be the descended map of \cref{lem:descent-criterion}.  For \(t\ge 0\),
	define
	\[
		\mathcal F_t^{(k)}
		:=
		\ker\bigl(\theta_k\circ \widehat\Phi_t^{(k)}\bigr)
		\subseteq S_k\times S_k,
	\]
	that is,
	\[
		(a,b)\in \mathcal F_t^{(k)}
		\iff
		\theta_k\bigl(\widehat\Phi_t^{(k)}(a)\bigr)
		=
		\theta_k\bigl(\widehat\Phi_t^{(k)}(b)\bigr).
	\]
\end{definition}
\begin{remark}[No independent finite-level observer]
	Finite-level observer data are not primitive.
	If \(\theta_\infty\) is \(k\)-level admissible, then \(\theta_k\)
	exists uniquely.  If \(\theta_\infty\) is not \(k\)-level admissible,
	no such \(\theta_k\) exists.  Any independent specification of a
	finite-level observer map would therefore constitute additional
	structure beyond the closed-system boundary.
\end{remark}
% ------------------------------------------------------------
\section{Determined interleaving across refinement}
% ------------------------------------------------------------
We now prove that observer-induced temporal identifications descend
along refinement whenever observer descent is available at adjacent
levels.
\begin{theorem}[Interleaving]
	\label{thm:interleaving}
	Assume that \(\theta_\infty\) is both \((k+1)\)-level admissible and
	\(k\)-level admissible.  Then for every \(t\ge 0\),
	\[
		(\phi_{k+1,k}\times \phi_{k+1,k})
		\bigl(\mathcal F_t^{(k+1)}\bigr)
		\subseteq
		\mathcal F_t^{(k)}.
	\]
\end{theorem}
\begin{proof}
	Let
	\[
		(u,v)\in \mathcal F_t^{(k+1)}.
	\]
	Then
	\[
		\theta_{k+1}\bigl(\widehat\Phi_t^{(k+1)}(u)\bigr)
		=
		\theta_{k+1}\bigl(\widehat\Phi_t^{(k+1)}(v)\bigr).
	\]
	Choose \(x,y\in S_\infty\) with
	\[
		p_{k+1}(x)=u,
		\qquad
		p_{k+1}(y)=v.
	\]
	Using
	\[
		\theta_{k+1}\circ p_{k+1}=\theta_\infty
	\]
	and
	\[
		p_{k+1}\circ \widehat\Phi_t^\infty
		=
		\widehat\Phi_t^{(k+1)}\circ p_{k+1},
	\]
	we obtain
	\[
		\theta_\infty\bigl(\widehat\Phi_t^\infty(x)\bigr)
		=
		\theta_\infty\bigl(\widehat\Phi_t^\infty(y)\bigr).
	\]
	Now use
	\[
		p_k=\phi_{k+1,k}\circ p_{k+1},
		\qquad
		p_k\circ \widehat\Phi_t^\infty
		=
		\widehat\Phi_t^{(k)}\circ p_k,
		\qquad
		\theta_k\circ p_k=\theta_\infty.
	\]
	Then
	\[
		\theta_k\bigl(\widehat\Phi_t^{(k)}(\phi_{k+1,k}(u))\bigr)
		=
		\theta_k\bigl(\widehat\Phi_t^{(k)}(\phi_{k+1,k}(v))\bigr),
	\]
	which is exactly the statement that
	\[
		(\phi_{k+1,k}(u),\phi_{k+1,k}(v))
		\in
		\mathcal F_t^{(k)}.
	\]
\end{proof}
\begin{corollary}[Refinement contraction]
	Define
	\[
		d_k(a,b):=\inf\{t\ge0:(a,b)\in \mathcal F_t^{(k)}\},
	\]
	whenever the infimum is meaningful.  Then
	\[
		d_k\bigl(\phi_{k+1,k}(u),\phi_{k+1,k}(v)\bigr)
		\le
		d_{k+1}(u,v).
	\]
\end{corollary}
\begin{proof}
	If
	\[
		(u,v)\in \mathcal F_t^{(k+1)},
	\]
	then \cref{thm:interleaving} gives
	\[
		(\phi_{k+1,k}(u),\phi_{k+1,k}(v))
		\in \mathcal F_t^{(k)}.
	\]
	Taking the infimum over all such \(t\) yields the claim.
\end{proof}
% ------------------------------------------------------------
\section{Spacetime as a joint two-parameter limit}
% ------------------------------------------------------------
\Cref{thm:Aut-functorial,thm:interleaving} produce two independent
directions of
structure.
\begin{itemize}
	\item
	      Refinement yields the inverse system
	      \[
		      \cdots \to S_{k+1}\to S_k\to\cdots .
	      \]
	\item
	      Observer-induced temporal coarse-graining yields, at each admissible
	      level \(k\), the quotient family
	      \[
		      S_k/\mathcal F_t^{(k)}.
	      \]
\end{itemize}
Consider therefore the two-parameter diagram
\[
	(k,t)\longmapsto S_k/\mathcal F_t^{(k)},
\]
restricted to those pairs \((k,t)\) for which \(\theta_\infty\)
descends to level \(k\) and for which the temporal bonding maps are
defined.
The refinement-direction bonding maps are supplied by
\cref{thm:interleaving}.  Temporal bonding maps exist on any subdomain
of \(t\) on which
\[
	t\longmapsto \mathcal F_t^{(k)}
\]
is monotone.
\begin{definition}[Spacetime object]
	The \emph{spacetime object} is
	\[
		\mathsf{ST}
		:=
		\varprojlim_{(k,t)}
		\bigl(S_k/\mathcal F_t^{(k)}\bigr),
	\]
	taken over the two-parameter domain on which the quotients and bonding
	maps are defined.
\end{definition}
\begin{remark}[Interpretation]
	The object \(\mathsf{ST}\) packages two distinct structural directions:
	\begin{enumerate}[label=\textup{(\roman*)}]
		\item
		      spatial refinement, represented by the inverse-system direction in
		      \(k\);
		\item
		      temporal coarsening, represented by the observer-induced quotient
		      direction in \(t\).
	\end{enumerate}
	It is therefore the first object in the stack in which the two limiting
	directions coexist in a single canonical construction.
\end{remark}
\begin{remark}[Determined versus observer-relative data]
	Up to this point, the following data are determined:
	\begin{itemize}
		\item
		      the intrinsic symmetry group
		      \[
			      G=\Aut(U,\mathcal C);
		      \]
		\item
		      comparison preservation of admissible dynamics under transport closure;
		\item
		      the functorial action of \(\Aut(U,\mathcal C)\) on the refinement
		      inverse system;
		\item
		      the induced action on the Stone limit;
		\item
		      the refinement-direction interleaving maps;
		\item
		      the corresponding two-parameter spacetime object on every domain where
		      observer descent and temporal monotonicity are available.
	\end{itemize}
	What remains observer-relative is the choice of reduction
	\[
		\theta_\infty\colon S_\infty\to Y,
	\]
	and hence the set of finite levels at which descent occurs, together
	with the temporal domain on which the induced filtration is monotone.
\end{remark}
\begin{remark}[Position in the stack]
	This chapter does not yet introduce curvature.
	Its output is the structural arena on which later curvature data will
	be realized:
	\[
		\begin{aligned}
			 & \text{transport-closed dynamics}
			\Longrightarrow
			\Aut(U,\mathcal C)\text{-action}
			\\
			 & \Longrightarrow
			\text{refinement/observer interleaving}
			\Longrightarrow
			\mathsf{ST}.
		\end{aligned}
	\]
	The curvature chapters begin only after this arena has been fixed.
\end{remark}

\section{Conclusion}
\label{sec:stitching-conclusion}
This chapter fixes the first joint refinement--time arena in the stack:
the two-parameter inverse-limit object \(\mathsf{ST}\), where intrinsic
refinement and observer-induced temporal coarsening are represented within
one coherent construction.

Accordingly, \cref{ch:christoffel} uses the stitched arena \(\mathsf{ST}\)
to realize the first intrinsic transport obstruction as smooth
curvature, while \cref{ch:causality,ch:connection} reuse that same arena
later in the causal and connection-first developments.